\documentclass[11pt,a4paper]{article}

\usepackage[utf8]{inputenc}
\usepackage[T1]{fontenc}
\usepackage[english]{babel}
\usepackage{lmodern}
\usepackage{microtype}
\usepackage{geometry}
\usepackage{amsmath,amssymb,amsthm,mathtools,bm}
\usepackage{booktabs}
\usepackage{array}
\usepackage{xcolor}
\usepackage{enumitem}
\usepackage{hyperref}

\geometry{margin=2.65cm}
\hypersetup{
  colorlinks=true,
  linkcolor=blue!55!black,
  citecolor=blue!55!black,
  urlcolor=blue!55!black,
  pdftitle={Pandrosion, Thales Scaling, and a Steffensen Acceleration},
  pdfauthor={Ivan Besevic}
}

\setlist[itemize]{leftmargin=1.8em,itemsep=0.25em,topsep=0.35em}
\setlist[enumerate]{leftmargin=1.8em,itemsep=0.25em,topsep=0.35em}

\theoremstyle{plain}
\newtheorem{theorem}{Theorem}[section]
\newtheorem{proposition}[theorem]{Proposition}
\newtheorem{lemma}[theorem]{Lemma}
\newtheorem{corollary}[theorem]{Corollary}

\theoremstyle{definition}
\newtheorem{definition}[theorem]{Definition}
\newtheorem{example}[theorem]{Example}

\theoremstyle{remark}
\newtheorem{remark}[theorem]{Remark}

\newcommand{\C}{\mathbb{C}}
\newcommand{\R}{\mathbb{R}}
\newcommand{\eps}{\varepsilon}
\newcommand{\dd}{\mathrm{d}}
\newcommand{\rootp}[2]{\sqrt[#1]{#2}}
\newcommand{\cbrt}[1]{\sqrt[3]{#1}}
\newcommand{\Sp}{S_p}
\newcommand{\hp}{h_{p,x}}
\newcommand{\stef}{\sigma_{p,x}}

\newenvironment{takeaway}{%
  \begin{center}\begin{minipage}{0.92\linewidth}\small\itshape\color{blue!45!black}%
}{%
  \end{minipage}\end{center}%
}

\title{\bfseries Pandrosion, Thales Scaling, and a Steffensen Acceleration\\[-0.1em]
\large A readable account of a geometric scheme for the equation $z^p=x$}
\author{Ivan \textsc{Besevic}\\[-0.2em]
\small Independent Mathematical Research\\[-0.2em]
\small \texttt{github.com/ivan-fr/poussiere\_de\_x}}
\date{April 2026}

\begin{document}
\maketitle

\begin{abstract}
Pandrosion of Alexandria is associated with a geometric approximation to the
problem of inserting proportional means, and hence to the classical problem of
extracting cube roots.  This paper rewrites that construction in modern
notation and follows the consequences as far as possible while keeping the
geometry visible.

The central object is a derivative-free fixed-point map for the equation
$z^p=x$.  On the real line it takes the universal form
\[
  h_{p,x}(s)=1-\frac{x-1}{x(1+s+\cdots+s^{p-1})}.
\]
Its fixed point is $s_*=x^{-1/p}$, and the desired root is recovered by
$v=x s_*^{p-1}=x^{1/p}$.  The contraction factor is explicit:
\[
  \lambda_{p,x}
  =\frac{(\alpha-1)\sum_{k=1}^{p-1}k\alpha^{p-1-k}}
         {\sum_{k=0}^{p-1}\alpha^k},
  \qquad \alpha=x^{1/p}.
\]
This expression explains why large targets slow the raw construction and why
homothetic preconditioning, or ``Thales scaling'', can dramatically reduce the
number of parallel constructions needed.

The map is then accelerated by Aitken--Steffensen's $\Delta^2$ transform,
which turns the linearly convergent Pandrosion iterator into a quadratically
convergent derivative-free scheme.  The complex extension is obtained by
working near the $p$ cyclotomic anchors of $x^{-1}$ and by seeding points into
local contraction disks.  The final section explains how the same fixed-point
idea extends from monomials to general monic polynomials.

The paper is intentionally modest in scope: it does not claim a universal
advance over modern polynomial solvers.  Its purpose is to present, in a clear
and self-contained form, the mathematics behind a concrete geometric
root-finding scheme, its residual profile, its scaling principle, and its
complex multi-start extension.
\end{abstract}

\begin{takeaway}
\textbf{One-sentence summary.}
Pandrosion supplies a geometric fixed-point map; Steffensen accelerates it;
homothety makes the geometry smaller before the map starts working.
\end{takeaway}

\tableofcontents
\newpage

%=====================================================================
\section{A reader's guide}
%=====================================================================

The paper is organized around one question:

\begin{quote}
\emph{What does Pandrosion's ruler-and-parallels construction become when it is
written as an iteration?}
\end{quote}

The answer is surprisingly compact.  A single rational map appears, its fixed
point is explicit, and its convergence factor can be computed in closed form.
The rest of the paper refines this basic observation.

\begin{itemize}
  \item Sections~\ref{sec:cube}--\ref{sec:general} derive the iteration from
  the cube-root case and then state the general $p$-th root form.
  \item Sections~\ref{sec:lambda}--\ref{sec:bounds} analyze the residual
  profile: the contraction factor, monotonicity, and error bounds.
  \item Sections~\ref{sec:start}--\ref{sec:scaling} discuss the practical
  improvements: a better first point, Steffensen acceleration, and scaling.
  \item Sections~\ref{sec:complex}--\ref{sec:lean} explain the complex
  multi-start picture and the formal-verification layer described in the
  original development.
  \item Section~\ref{sec:general-polynomials} explains what remains true for
  a general monic polynomial.
\end{itemize}

The exposition is deliberately less compressed than the original manuscript.
The goal is not to hide the algebra, but to make clear what each formula is
for.

%=====================================================================
\section{The cube-root construction}
\label{sec:cube}
%=====================================================================

Pandrosion's construction begins with the ancient problem of doubling the cube:
constructing a segment of length $\cbrt{2}$ from a unit segment.  Algebraically,
this is the problem of inserting two proportional means $m_1,m_2$ between $1$
and $2$:
\[
  \frac{1}{m_1}=\frac{m_1}{m_2}=\frac{m_2}{2}.
\]
The solution is $m_1=\cbrt{2}$ and $m_2=\cbrt{4}$.

Pandrosion's geometric procedure does not produce the exact value in finitely
many steps.  It produces a sequence of approximations.  This is the key modern
viewpoint: the construction is an iteration, and the error sequence of that
iteration is its residual profile.

In the classical cube-root case one may write the construction with an internal
length $u_n$ and a visible readout $v_n$:
\begin{equation}
  u_{n+1}
  =\frac{u_n}{2-2\left(\frac{4-u_n}{4}\right)^3},
  \qquad
  v_n=2\left(\frac{4-u_n}{4}\right)^2.
  \label{eq:classical-u}
\end{equation}
The number $u_n$ belongs to the diagram.  The number $v_n$ is the approximation
that is read from the diagram.

The substitution
\[
  s_n=\frac{4-u_n}{4}
\]
removes the distracting constants.  Since $u_n=4(1-s_n)$, the readout becomes
$v_n=2s_n^2$, and the iteration becomes
\begin{equation}
  s_{n+1}
  =\frac{2s_n^2+2s_n+1}{2s_n^2+2s_n+2}
  =1-\frac{1}{2(1+s_n+s_n^2)}.
  \label{eq:cube-s}
\end{equation}

\begin{proposition}[Cube-root fixed point]
The map in~\eqref{eq:cube-s} has the fixed point
\[
  s_*=2^{-1/3}.
\]
Consequently the readout satisfies
\[
  v_n=2s_n^2 \longrightarrow 2s_*^2=\cbrt{2}.
\]
\end{proposition}

\begin{proof}
The fixed-point equation is
\[
  s=1-\frac{1}{2(1+s+s^2)}.
\]
Multiplying by $2(1+s+s^2)$ gives
\[
  2s(1+s+s^2)=2(1+s+s^2)-1,
\]
which simplifies to $2s^3=1$.  Hence $s=2^{-1/3}$ in the positive real branch.
The readout identity follows immediately.
\end{proof}

\begin{takeaway}
The diagram computes $\cbrt{2}$ indirectly.  It first converges to
$2^{-1/3}$, and the visible root is recovered by the readout $v=2s^2$.
\end{takeaway}

%=====================================================================
\section{Why the polynomial \texorpdfstring{$S_p$}{Sp} appears}
\label{sec:geometry}
%=====================================================================

The polynomial
\[
  \Sp(s)=1+s+\cdots+s^{p-1}
\]
appears for a geometric reason.  Thales' theorem constructs proportional
segments.  Repeating the construction builds a finite geometric progression of
ratios.  Algebraically that progression is exactly $\Sp$.

The elementary identity behind the construction is
\begin{equation}
  1-s^p=(1-s)(1+s+\cdots+s^{p-1})=(1-s)\Sp(s).
  \label{eq:telescopic-Sp}
\end{equation}
Thus, whenever a diagram compares $1$ and $s^p$ through successive proportional
segments, the denominator $\Sp(s)$ is the natural object that records the
accumulated proportional means.

This identity is the bridge between the geometric construction and the
fixed-point map.

%=====================================================================
\section{The universal real-line iteration}
\label{sec:general}
%=====================================================================

Let $p\ge 2$ and $x>0$.  Define
\begin{equation}
  \Sp(s)=\sum_{k=0}^{p-1}s^k,
  \qquad
  \hp(s)=1-\frac{x-1}{x\Sp(s)}.
  \label{eq:hp-def}
\end{equation}
The Pandrosion iteration is
\begin{equation}
  s_{n+1}=\hp(s_n).
  \label{eq:pandrosion-iteration}
\end{equation}
The desired root is not $s_n$ itself.  It is read from $s_n$ by
\begin{equation}
  v_n=x s_n^{p-1}.
  \label{eq:readout}
\end{equation}

\begin{theorem}[Fixed point and readout]
\label{thm:fixed-point}
Let $x>1$ and $p\ge 2$.  The map $\hp$ has the positive fixed point
\[
  s_*=x^{-1/p}.
\]
At this fixed point the readout is
\[
  x s_*^{p-1}=x^{1/p}.
\]
\end{theorem}

\begin{proof}
The fixed-point equation $s=\hp(s)$ is
\[
  1-s=\frac{x-1}{x\Sp(s)}.
\]
Using~\eqref{eq:telescopic-Sp}, this is equivalent to
\[
  x(1-s)\Sp(s)=x-1,
  \qquad\text{hence}\qquad
  x(1-s^p)=x-1.
\]
Therefore $xs^p=1$, so $s=x^{-1/p}$ in the positive branch.  The readout is
then $x(x^{-1/p})^{p-1}=x^{1/p}$.
\end{proof}

\begin{remark}
The map $\hp$ is derivative-free in its definition.  It is built from the
finite geometric sum $\Sp$, not from the derivative of $s^p-x$.
\end{remark}

%=====================================================================
\section{The contraction ratio}
\label{sec:lambda}
%=====================================================================

The speed of the raw Pandrosion iteration is governed by the derivative of
$\hp$ at the fixed point.  This derivative is not used to define the method; it
is used only to analyze the residual profile.

Let
\[
  \alpha=x^{1/p}.
\]
Then $s_*=\alpha^{-1}$.

\begin{theorem}[Closed-form contraction factor]
\label{thm:lambda}
For $x>1$ and $p\ge2$, the linear contraction factor of the Pandrosion
iteration at the fixed point is
\begin{equation}
  \lambda_{p,x}=\hp'(s_*)
  =\frac{(\alpha-1)\sum_{k=1}^{p-1}k\alpha^{p-1-k}}
         {\sum_{k=0}^{p-1}\alpha^k}.
  \label{eq:lambda}
\end{equation}
In particular,
\[
  \lambda_{2,x}=\frac{\alpha-1}{\alpha+1},
  \qquad
  \lambda_{3,x}=\frac{(\alpha-1)(\alpha+2)}{\alpha^2+\alpha+1}.
\]
\end{theorem}

\begin{proof}
Differentiate~\eqref{eq:hp-def}:
\[
  \hp'(s)=\frac{x-1}{x}\frac{\Sp'(s)}{\Sp(s)^2}.
\]
Substitute $s=\alpha^{-1}$ and use
\[
  \Sp(\alpha^{-1})
  =\sum_{k=0}^{p-1}\alpha^{-k}
  =\alpha^{1-p}\sum_{k=0}^{p-1}\alpha^k.
\]
The same change of powers in $\Sp'(\alpha^{-1})$ yields~\eqref{eq:lambda}.
\end{proof}

The formula immediately explains the qualitative behavior:

\begin{itemize}
  \item if $x$ is close to $1$, then $\lambda_{p,x}$ is small and the raw
  geometry converges quickly;
  \item if $x$ is large, then $\lambda_{p,x}$ approaches $1$ and the raw
  geometry converges slowly;
  \item increasing $p$ at fixed large target makes the unscaled construction
  more delicate.
\end{itemize}

\begin{proposition}[Large-$p$ limit at fixed target]
For fixed $x>1$,
\[
  \lim_{p\to\infty}\lambda_{p,x}
  =1-\frac{\log x}{x-1}.
\]
\end{proposition}

\begin{proof}[Proof sketch]
Write $\alpha=x^{1/p}=e^{(\log x)/p}$ and replace the sums in
\eqref{eq:lambda} by Riemann sums on $[0,1]$.  The resulting quotient reduces
to the displayed expression.
\end{proof}

\begin{takeaway}
The number $\lambda_{p,x}$ is the fingerprint of Pandrosion's residual
profile.  Scaling works because it changes $x$ before this fingerprint is
formed.
\end{takeaway}

%=====================================================================
\section{Error bounds and iteration counts}
\label{sec:bounds}
%=====================================================================

The local statement is the usual one for a contracting fixed-point map.  If an
interval $I$ contains $s_*$ and
\[
  \sup_{s\in I}|\hp'(s)|\le \Lambda<1,
\]
then every orbit that remains in $I$ satisfies
\begin{equation}
  |s_n-s_*|\le \Lambda^n |s_0-s_*|.
  \label{eq:basic-error}
\end{equation}
Consequently the readout error is controlled by the same geometric decay,
with an additional Lipschitz constant for $s\mapsto xs^{p-1}$ on $I$.

A practical iteration budget follows from~\eqref{eq:basic-error}.  To reach
$|s_n-s_*|\le \tau$, it is enough to take
\begin{equation}
  n\ge
  \frac{\log(\tau/|s_0-s_*|)}{\log \Lambda}.
  \label{eq:budget}
\end{equation}
The denominator is negative, so a smaller $\Lambda$ means a smaller budget.
This is the numerical meaning of the geometric statement: fewer effective
Thales parallels are needed.

%=====================================================================
\section{The case \texorpdfstring{$0<x<1$}{0 < x < 1}}
\label{sec:x-less-one}
%=====================================================================

When $0<x<1$, the fixed point $s_*=x^{-1/p}$ lies outside $(0,1)$.  The direct
real-line diagram that was natural for $x>1$ therefore no longer lives in the
same interval.  One can still analyze the rational map, but the clean geometric
contraction picture changes.

The practical resolution is simple: use the reciprocal target.  Since
\[
  x^{1/p}=\frac{1}{(1/x)^{1/p}},
\]
one applies the $x>1$ theory to $1/x$ and then inverts the result.  This is the
first appearance of a principle that becomes important in the complex theory:
changing charts is often better than forcing one chart to handle every scale.

%=====================================================================
\section{A better starting point}
\label{sec:start}
%=====================================================================

The raw iteration needs an initial point.  The most natural choice is not
always the fastest.  A simple geometric improvement is to start from the image
of the neutral point $1$:
\begin{equation}
  s_0^{\mathrm{opt}}=\hp(1)=1-\frac{x-1}{xp}.
  \label{eq:startopt}
\end{equation}
This point already incorporates one full proportional-mean step.  In the
language of the diagram, one avoids spending future parallels to do work that
can be done once at initialization.

\begin{proposition}[Canonical start]
The start $s_0^{\mathrm{opt}}=\hp(1)$ is the canonical one-step start produced
by the same Pandrosion geometry.  Near $x=1$ it satisfies
\[
  s_0^{\mathrm{opt}}=x^{-1/p}+O((x-1)^2).
\]
\end{proposition}

\begin{proof}
The first statement is the definition.  The expansion follows from
$x^{-1/p}=1-(x-1)/p+O((x-1)^2)$ and from
$1-(x-1)/(xp)=1-(x-1)/p+O((x-1)^2)$.
\end{proof}

\begin{takeaway}
Starting-point optimization is not an extra algorithmic layer.  It is the first
Pandrosion step used deliberately, before the iteration budget begins.
\end{takeaway}

%=====================================================================
\section{Steffensen acceleration}
\label{sec:steffensen}
%=====================================================================

The raw Pandrosion map is geometrically meaningful but only linearly
convergent.  Aitken--Steffensen acceleration turns any sufficiently regular
linearly convergent fixed-point iteration into a quadratically convergent one,
without requiring derivatives.

For $h=\hp$, define
\begin{equation}
  \stef(s)
  =s-\frac{(h(s)-s)^2}{h(h(s))-2h(s)+s}.
  \label{eq:steffensen}
\end{equation}

\begin{theorem}[Quadratic acceleration]
Assume $h$ is analytic near $s_*$, $h(s_*)=s_*$, and $h'(s_*)\ne 1$.  Then the
Steffensen map $\stef$ has the same fixed point and satisfies
\[
  \stef(s)-s_*=C(s-s_*)^2+O((s-s_*)^3)
\]
for an explicit constant $C$ depending on the first derivatives of $h$ at
$s_*$.  Thus the accelerated Pandrosion iteration is locally quadratic.
\end{theorem}

\begin{proof}[Proof sketch]
Expand $h(s_*+e)$ as
\[
  s_*+\lambda e+a e^2+O(e^3).
\]
Substitute this expansion into~\eqref{eq:steffensen}.  The linear term cancels,
leaving a quadratic leading term.
\end{proof}

Steffensen acceleration is especially natural here.  The underlying geometric
map remains derivative-free, and the acceleration uses only evaluations of the
same map.

%=====================================================================
\section{Thales scaling}
\label{sec:scaling}
%=====================================================================

Scaling is the most important practical improvement.  Suppose $x$ is large.
Instead of applying Pandrosion directly to $x$, write
\begin{equation}
  x=A x',
  \label{eq:scaling-factor}
\end{equation}
where $A$ is chosen so that $x'$ is close to $1$.  Then
\[
  x^{1/p}=A^{1/p}(x')^{1/p}.
\]
One computes $(x')^{1/p}$ by Pandrosion and multiplies back by $A^{1/p}$.

The point is not cosmetic.  The contraction factor is controlled by $x'$, not
by $x$.  Since $\lambda_{p,x}$ becomes small when $x$ is close to $1$, the
scaled problem may require far fewer effective Thales constructions.

\begin{takeaway}
Homothety does not merely normalize the input.  It changes the geometry before
Pandrosion begins.  A good homothety reduces the number of proportional layers
the construction must build.
\end{takeaway}

This is the one-dimensional ancestor of the geometric scaling used later in
multivariate Pandrosion solvers: first move the problem into a stable chart,
then apply the local correction.

%=====================================================================
\section{Comparison with standard methods}
\label{sec:comparison}
%=====================================================================

For $f(u)=u^p-x$, Newton's method is
\[
  u_{n+1}=u_n-\frac{u_n^p-x}{p u_n^{p-1}}.
\]
It is quadratically convergent near a simple root, but it uses the derivative
of $f$.  Halley's method and the higher Schröder--König family can converge
with even higher order.

Pandrosion is different in spirit.  Its raw form is linearly convergent, but it
arises from a proportional geometry and uses only the finite sum $\Sp$.  After
Steffensen acceleration, it becomes derivative-free and quadratic.  After
scaling, it can be made numerically competitive in regimes where the geometric
preconditioner makes $x'$ close to $1$.

\begin{center}
\begin{tabular}{@{}lll@{}}
\toprule
Method & Local order & Uses derivative? \\
\midrule
Raw Pandrosion & linear & no \\
Steffensen--Pandrosion & quadratic & no \\
Newton & quadratic & yes \\
Halley & cubic & yes \\
\bottomrule
\end{tabular}
\end{center}

The intended message is not that Pandrosion dominates Newton.  It usually does
not.  The message is that Pandrosion supplies a geometric fixed-point operator
whose residual profile is exactly analyzable and whose performance improves
substantially under scaling and Steffensen acceleration.

%=====================================================================
\section{Complex extension}
\label{sec:complex}
%=====================================================================

The real-line analysis identifies one positive fixed point.  Over $\C$, the
equation $s^p=1/x$ has $p$ fixed points:
\begin{equation}
  \gamma_k=x^{-1/p}\zeta_p^k,
  \qquad
  \zeta_p=e^{2\pi i/p},
  \qquad k=0,\dots,p-1.
  \label{eq:anchors}
\end{equation}
These are the cyclotomic anchors of the complex Pandrosion map.

The complex dynamics has non-trivial basin structure.  A practical way to avoid
this difficulty is to seed points directly into a known local contraction disk
around the desired anchor.

\begin{definition}[Targeted seed]
Given an anchor $\gamma$ and a parameter $0<\eps<1$, define
\[
  \operatorname{seed}_{\gamma,\eps}(z)
  =\gamma+\eps(z-\gamma).
\]
\end{definition}

The seed is a homothety centered at $\gamma$.  It pulls any point $z$ closer to
the target anchor.  For $\eps$ small enough, the seeded point lies in a local
quadratic-convergence region for the Steffensen--Pandrosion map.

\begin{theorem}[Seeded complex convergence]
Let $\gamma$ be one of the anchors~\eqref{eq:anchors}.  There exists a radius
$r_\gamma>0$ such that if
\[
  |\operatorname{seed}_{\gamma,\eps}(z)-\gamma|<r_\gamma,
\]
then the Steffensen--Pandrosion orbit converges quadratically to $\gamma$.
Moreover, choosing $\eps$ small enough enforces this condition for every fixed
$z\in\C$.
\end{theorem}

\begin{proof}[Proof sketch]
The Steffensen map is analytic near a simple fixed point and has a quadratic
fixed point there.  Hence a sufficiently small disk around $\gamma$ is mapped
into itself with quadratic contraction.  The seed map shrinks the distance to
$\gamma$ by the factor $\eps$.
\end{proof}

Once $\gamma_k$ is obtained, the corresponding root of $x$ is reconstructed by
\[
  \rho_k=x\gamma_k^{p-1}=x^{1/p}\zeta_p^{-k}.
\]
Thus the complex scheme is a multi-start method indexed by the cyclotomic
anchors.

%=====================================================================
\section{Measurability and continuity of the solution selector}
\label{sec:measurability}
%=====================================================================

A target selector chooses which anchor should receive a given starting point.
The simplest selector is the nearest-anchor rule: assign $z$ to the anchor
$\gamma_k$ that minimizes $|z-\gamma_k|$.

This selector is discontinuous on the Voronoi boundaries between anchors, but
these boundaries have Lebesgue measure zero in the plane.  Away from them, the
chosen anchor is locally constant, and the seeded iteration depends
continuously on the input.

\begin{proposition}[A.e. continuity]
The solution map induced by the nearest-anchor selector and the seeded
Steffensen--Pandrosion iteration is Borel measurable everywhere and continuous
outside the Voronoi boundary of the anchor set.
\end{proposition}

\begin{proof}[Proof sketch]
Voronoi cells are defined by finitely many inequalities of continuous
functions, hence are Borel.  Their boundaries are finite unions of real
codimension-one sets.  Inside a cell, the same anchor is selected and the
iteration is a locally uniform limit of continuous functions.
\end{proof}

%=====================================================================
\section{Formal verification layer}
\label{sec:lean}
%=====================================================================

The original development reports a Lean~4 / Mathlib v4.7.0 corpus formalizing
the elementary algebraic backbone of the construction.  Its role is not to
certify every numerical experiment.  Its role is to pin down the identities on
which the scheme rests:

\begin{itemize}
  \item the finite-sum identity $(1-s)\Sp(s)=1-s^p$;
  \item the fixed-point equation for $\hp$;
  \item the reconstruction identity $x s_*^{p-1}=x^{1/p}$;
  \item the compatibility of seeding, target selection, and reconstruction in
  the complex multi-start setting.
\end{itemize}

The reported corpus contains 121 modules with no remaining \texttt{sorry} and
with axioms restricted to the standard Lean whitelist described in the original
project notes.  This is best understood as a formal audit of the algebraic
core, not as a claim that the numerical method is globally optimal.

%=====================================================================
\section{Beyond monomials}
\label{sec:general-polynomials}
%=====================================================================

The monomial equation $z^p=x$ is the clean case because the finite geometric
sum $\Sp$ gives a closed-form denominator.  Nevertheless the same fixed-point
idea extends to any monic polynomial
\[
  P(z)=z^p+R(z).
\]
Define
\begin{equation}
  h_P(s)=1-\frac{R(s)+1}{\Sp(s)}.
  \label{eq:general-hp}
\end{equation}

\begin{proposition}[Fixed points of the generalized Pandrosion map]
A point $s$ is a fixed point of $h_P$ if and only if it is a root of $P$.
\end{proposition}

\begin{proof}
The equation $s=h_P(s)$ is
\[
  1-s=\frac{R(s)+1}{\Sp(s)}.
\]
Multiplying by $\Sp(s)$ and using $(1-s)\Sp(s)=1-s^p$ gives
\[
  1-s^p=R(s)+1,
\]
which is equivalent to $s^p+R(s)=0$.
\end{proof}

This extension preserves the derivative-free character of the fixed-point map,
but it loses the clean closed-form contraction profile of the monomial case.
The derivative $h_P'(\alpha)$ depends on the full polynomial $R$, and roots may
be near singularities of the rational map.  Therefore the monomial theory is
not automatically a uniform theory for all polynomials.

\begin{takeaway}
The monomial case is special because geometry, algebra, and convergence all
collapse to one finite sum.  General polynomials keep the fixed-point idea but
not the same closed-form speed guarantees.
\end{takeaway}

%=====================================================================
\section{Conclusion}
\label{sec:conclusion}
%=====================================================================

Pandrosion's construction can be read as a root-finding algorithm.  In modern
notation, it is the fixed-point scheme
\[
  s_{n+1}=1-\frac{x-1}{x(1+s_n+\cdots+s_n^{p-1})}.
\]
Its fixed point is $x^{-1/p}$, its readout is $x^{1/p}$, and its contraction
ratio is the explicit quantity~\eqref{eq:lambda}.

Three ideas make the method practically interesting:

\begin{enumerate}
  \item \textbf{Starting-point optimization:} begin from a point already shaped
  by the Pandrosion geometry.
  \item \textbf{Steffensen acceleration:} upgrade the derivative-free fixed
  point map from linear to quadratic convergence.
  \item \textbf{Thales scaling:} use homothety to reduce the target before the
  geometric construction starts.
\end{enumerate}

The complex extension adds cyclotomic anchors and seeded local convergence.
The polynomial extension shows that the fixed-point idea is not restricted to
monomials, although the monomial case is the one where the residual profile is
fully explicit.

The main lesson is geometric: before asking an iteration to converge, one may
change the space in which it works.  In the Pandrosion viewpoint, homothety is
not a numerical trick.  It is part of the construction itself.

%=====================================================================
\begin{thebibliography}{10}

\bibitem{pappus}
Pappus of Alexandria,
\textit{Collectio}, Book III, ca. 340 AD.
Critical edition: F. Hultsch, \textit{Pappi Alexandrini Collectionis quae supersunt},
Weidmann, Berlin, 1876--1878.

\bibitem{mactutor_pandrosion}
J. J. O'Connor and E. F. Robertson,
``Pandrosion of Alexandria,''
\textit{MacTutor History of Mathematics}, University of St Andrews, 2020.

\bibitem{mactutor_cube}
J. J. O'Connor and E. F. Robertson,
``Doubling the cube,''
\textit{MacTutor History of Mathematics}, University of St Andrews, 2002.

\bibitem{steffensen}
J. F. Steffensen,
``Remarks on iteration,''
\textit{Skandinavisk Aktuarietidskrift} \textbf{16} (1933), 64--72.

\bibitem{aitken}
A. C. Aitken,
``On Bernoulli's numerical solution of algebraic equations,''
\textit{Proceedings of the Royal Society of Edinburgh} \textbf{46} (1926), 289--305.

\bibitem{ostrowski}
A. M. Ostrowski,
\textit{Solution of Equations and Systems of Equations},
2nd ed., Academic Press, 1966.

\bibitem{traub}
J. F. Traub,
\textit{Iterative Methods for the Solution of Equations},
Prentice-Hall, 1964.

\bibitem{householder}
A. S. Householder,
\textit{The Numerical Treatment of a Single Nonlinear Equation},
McGraw-Hill, 1970.

\bibitem{leanprover}
L. de Moura and S. Ullrich,
``The Lean 4 theorem prover and programming language,''
CADE-28, Lecture Notes in Computer Science, Springer, 2021.

\bibitem{mathlib}
The Mathlib Community,
\textit{Mathlib: A unified library of mathematics formalized}.
\url{https://github.com/leanprover-community/mathlib4}

\end{thebibliography}

\end{document}
